\documentclass[preview]{standalone}

\usepackage{amsmath}
\usepackage{amssymb}
\usepackage{stellar}
\usepackage{definitions}
\usepackage{bettelini}

\begin{document}

\id{linearalgebra-gaussian-elimination}
\genpage

\section{Equivalent systems}

\begin{snippetdefinition}{linear-system-reduction-definition}{Reduction of a linear system}
    Let \(S\) be a linear system. A \emph{reduction} of \(S\) is a
    transformation of \(S\) into an equivalent linear system that is simpler
    to solve.
\end{snippetdefinition}

\begin{snippetdefinition}{equivalent-linear-systems-definition}{Equivalent linear systems}
    Two linear systems in the same unknowns are \emph{equivalent} if they have
    the same solution \set.
\end{snippetdefinition}

\begin{snippetlemma}{linear-equation-addition-operation}{Adding a multiple of one equation to another}
    Let \(S\) be a linear system over a \field \(\mathbb{F}\), and let
    \(E_i\) and \(E_j\) be two equations of \(S\), with \(i\neq j\).
    If \(\lambda\in\mathbb{F}\), then replacing \(E_i\) by
    \[
        E_i+\lambda E_j
    \]
    produces a linear system equivalent to \(S\).
\end{snippetlemma}

\begin{snippetproof}{linear-equation-addition-operation-proof}{linear-equation-addition-operation}{Adding a multiple of one equation to another}
    Let \(\overline{x}\) be a solution of \(S\). Then
    \(\overline{x}\) satisfies both \(E_i\) and \(E_j\). If
    \[
        E_i:\quad a_{i1}x_1+\cdots+a_{in}x_n=b_i
    \]
    and
    \[
        E_j:\quad a_{j1}x_1+\cdots+a_{jn}x_n=b_j,
    \]
    then the new equation is
    \[
        (a_{i1}+\lambda a_{j1})x_1+\cdots
        +(a_{in}+\lambda a_{jn})x_n
        =
        b_i+\lambda b_j.
    \]
    Substituting \(\overline{x}\) gives \(b_i+\lambda b_j\), so
    \(\overline{x}\) also satisfies the new equation.

    Conversely, the operation is reversible: from
    \(E_i'=E_i+\lambda E_j\) one recovers \(E_i=E_i'-\lambda E_j\). Hence
    the solution \set is unchanged.
\end{snippetproof}

\begin{snippetlemma}{linear-equation-swap-operation}{Swapping equations}
    Let \(S\) be a linear system. Swapping two equations of \(S\) produces a
    linear system equivalent to \(S\).
\end{snippetlemma}

\begin{snippetproof}{linear-equation-swap-operation-proof}{linear-equation-swap-operation}{Swapping equations}
    Swapping two equations changes only the order in which the conditions are
    written. A tuple of unknowns satisfies all equations before the swap if and
    only if it satisfies all equations after the swap.
\end{snippetproof}

\section{Gaussian elimination}

\begin{snippetdefinition}{gaussian-elimination-pivot-definition}{Pivot}
    During Gaussian elimination, a \emph{pivot} is a nonzero coefficient used
    to eliminate the entries below it in the same column.
\end{snippetdefinition}

\begin{snippetdefinition}{forward-elimination-definition}{Forward elimination}
    Let \(S\) be a linear system. \emph{Forward elimination} is the phase of
    Gaussian elimination in which pivots are used to eliminate unknowns from
    later equations. In the nonsingular square case, this phase transforms the
    system into a triangular system.
\end{snippetdefinition}

\begin{snippetdefinition}{back-substitution-definition}{Back substitution}
    Let \(S\) be a triangular linear system. \emph{Back substitution} is the
    process of solving \(S\) starting from the last equation and then moving
    upward through the preceding equations.
\end{snippetdefinition}

\begin{snippetexample}{gaussian-elimination-first-example}{Gaussian elimination}
    Consider the linear system
    \[
        \begin{cases}
            2u+v+w=5,\\
            4u-6v=-2,\\
            -2u+7v+2w=9.
        \end{cases}
    \]
    Subtracting \(2\) times the first equation from the second, and adding
    the first equation to the third, gives the equivalent system
    \[
        \begin{cases}
            2u+v+w=5,\\
            -8v-2w=-12,\\
            8v+3w=14.
        \end{cases}
    \]
    Adding the second equation to the third gives
    \[
        \begin{cases}
            2u+v+w=5,\\
            -8v-2w=-12,\\
            w=2.
        \end{cases}
    \]
    Back substitution gives
    \[
        w=2,
        \qquad
        v=1,
        \qquad
        u=1.
    \]
    Hence the unique solution is
    \[
        (u,v,w)=(1,1,2).
    \]
\end{snippetexample}

\begin{snippet}{gaussian-elimination-matrix-form}
    The system in the previous example has augmented \matrix
    \[
        \left[
            \begin{array}{ccc|c}
                2 & 1 & 1 & 5\\
                4 & -6 & 0 & -2\\
                -2 & 7 & 2 & 9
            \end{array}
        \right].
    \]
    The row operations
    \[
        R_2\leftarrow R_2-2R_1,
        \qquad
        R_3\leftarrow R_3+R_1
    \]
    give
    \[
        \left[
            \begin{array}{ccc|c}
                2 & 1 & 1 & 5\\
                0 & -8 & -2 & -12\\
                0 & 8 & 3 & 14
            \end{array}
        \right].
    \]
    Then
    \[
        R_3\leftarrow R_3+R_2
    \]
    gives
    \[
        \left[
            \begin{array}{ccc|c}
                2 & 1 & 1 & 5\\
                0 & -8 & -2 & -12\\
                0 & 0 & 1 & 2
            \end{array}
        \right].
    \]
    The pivots are \(2\), \(-8\), and \(1\).
\end{snippet}

\begin{snippetproposition}{gaussian-elimination-square-nonsingular-method}{Gaussian elimination for nonsingular square systems}
    Let \(S\) be a square linear system with \(n\) equations and \(n\)
    unknowns over a \field \(\mathbb{F}\). If forward elimination produces
    \(n\) pivots, then \(S\) has a unique solution, obtained by back
    substitution.
\end{snippetproposition}

\begin{snippetproof}{gaussian-elimination-square-nonsingular-method-proof}{gaussian-elimination-square-nonsingular-method}{Gaussian elimination for nonsingular square systems}
    Forward elimination uses reversible elementary equation operations, so the
    solution \set is unchanged. If \(n\) pivots are produced, the resulting
    system is triangular with a nonzero coefficient in each pivot position.
    The last equation determines the last unknown uniquely. Substituting this
    value into the previous equation determines the previous unknown uniquely.
    Continuing upward determines all unknowns uniquely.
\end{snippetproof}

\section{Singular cases}

\begin{snippetexample}{gaussian-elimination-inconsistent-zero-row-example}{Inconsistent zero row}
    Consider the system
    \[
        \begin{cases}
            2x-4y=7,\\
            x-2y=5.
        \end{cases}
    \]
    Its augmented \matrix is
    \[
        \left[
            \begin{array}{cc|c}
                2 & -4 & 7\\
                1 & -2 & 5
            \end{array}
        \right].
    \]
    After the operation
    \[
        R_2\leftarrow R_2-\frac{1}{2}R_1,
    \]
    one obtains
    \[
        \left[
            \begin{array}{cc|c}
                2 & -4 & 7\\
                0 & 0 & \frac{3}{2}
            \end{array}
        \right].
    \]
    The second row represents \(0=\frac{3}{2}\), which is impossible.
    Therefore the system has no solutions.
\end{snippetexample}

\begin{snippetexample}{gaussian-elimination-free-variable-example}{Free variable after elimination}
    Consider the system
    \[
        \begin{cases}
            2x-4y=6,\\
            x-2y=3.
        \end{cases}
    \]
    After the row operation
    \[
        R_2\leftarrow R_2-\frac{1}{2}R_1
    \]
    the augmented \matrix becomes
    \[
        \left[
            \begin{array}{cc|c}
                2 & -4 & 6\\
                0 & 0 & 0
            \end{array}
        \right].
    \]
    The second row represents the identity \(0=0\). The system is equivalent
    to \(x-2y=3\). If \(y=t\in\realnumbers\), then \(x=3+2t\). Thus the
    solution \set is
    \[
        \{(3+2t,t)\suchthat t\in\realnumbers\}.
    \]
\end{snippetexample}

\begin{snippetnote}{gaussian-elimination-zero-pivot-note}{Zero pivots}
    In a square linear system, the failure to produce a pivot in every column
    does not by itself decide whether the system has solutions. The augmented
    column determines whether the system is inconsistent or has free
    variables.
\end{snippetnote}

\begin{snippetexample}{gaussian-elimination-row-swap-example}{Using a row swap}
    Consider a system whose augmented \matrix becomes
    \[
        \left[
            \begin{array}{ccc|c}
                1 & 1 & 1 & b_1\\
                0 & 0 & 3 & b_2-2b_1\\
                0 & 2 & 4 & b_3-4b_1
            \end{array}
        \right]
    \]
    during elimination. The entry in the second column of the second row is
    zero, but swapping the second and third rows gives
    \[
        \left[
            \begin{array}{ccc|c}
                1 & 1 & 1 & b_1\\
                0 & 2 & 4 & b_3-4b_1\\
                0 & 0 & 3 & b_2-2b_1
            \end{array}
        \right].
    \]
    This \matrix has three pivots \(1\), \(2\), and \(3\). Hence the
    corresponding square system has a unique solution for every choice of
    \(b_1,b_2,b_3\).
\end{snippetexample}

\begin{snippetexample}{gaussian-elimination-singular-square-system-example}{A singular square system}
    Consider the square system
    \[
        \begin{cases}
            u+v+w=b_1,\\
            2u+2v+5w=b_2,\\
            4u+4v+8w=b_3.
        \end{cases}
    \]
    After eliminating the first column one obtains
    \[
        \left[
            \begin{array}{ccc|c}
                1 & 1 & 1 & b_1\\
                0 & 0 & 3 & b_2-2b_1\\
                0 & 0 & 4 & b_3-4b_1
            \end{array}
        \right].
    \]
    No pivot can be produced in the second column. Depending on the values of
    \(b_1,b_2,b_3\), the system may have infinitely many solutions or no
    solutions.
\end{snippetexample}

\begin{snippetexample}{gaussian-elimination-singular-infinite-solutions-example}{Infinite solutions in a singular case}
    In the singular system
    \[
        \begin{cases}
            u+v+w=b_1,\\
            2u+2v+5w=b_2,\\
            4u+4v+8w=b_3,
        \end{cases}
    \]
    take \(b_1=0\), \(b_2=3\), and \(b_3=4\). The reduced system is
    \[
        \begin{cases}
            u+v+w=0,\\
            3w=3,\\
            4w=4.
        \end{cases}
    \]
    Hence \(w=1\) and \(u+v+1=0\). If \(v=t\in\realnumbers\), then
    \(u=-1-t\), and the solution \set is
    \[
        \{(-1-t,t,1)\suchthat t\in\realnumbers\}.
    \]
\end{snippetexample}

\begin{snippetexample}{gaussian-elimination-singular-no-solutions-example}{No solutions in a singular case}
    In the singular system
    \[
        \begin{cases}
            u+v+w=b_1,\\
            2u+2v+5w=b_2,\\
            4u+4v+8w=b_3,
        \end{cases}
    \]
    take \(b_1=0\), \(b_2=4\), and \(b_3=3\). The last two reduced equations
    become
    \[
        3w=4,
        \qquad
        4w=3.
    \]
    These equations are incompatible, so the system has no solutions.
\end{snippetexample}

\section{Complete example}

\begin{snippetexample}{gaussian-elimination-four-unknowns-example}{Gaussian elimination in four unknowns}
    Consider the system
    \[
        \begin{cases}
            w+2x-y+z=6,\\
            -w+x+2y-z=3,\\
            2w-x+2y+2z=14,\\
            w+x-y+2z=8.
        \end{cases}
    \]
    Its augmented \matrix is
    \[
        \left[
            \begin{array}{rrrr|r}
                1 & 2 & -1 & 1 & 6\\
                -1 & 1 & 2 & -1 & 3\\
                2 & -1 & 2 & 2 & 14\\
                1 & 1 & -1 & 2 & 8
            \end{array}
        \right].
    \]
    After the operations
    \[
        R_2\leftarrow R_2+R_1,
        \quad
        R_3\leftarrow R_3-2R_1,
        \quad
        R_4\leftarrow R_4-R_1,
    \]
    one gets
    \[
        \left[
            \begin{array}{rrrr|r}
                1 & 2 & -1 & 1 & 6\\
                0 & 3 & 1 & 0 & 9\\
                0 & -5 & 4 & 0 & 2\\
                0 & -1 & 0 & 1 & 2
            \end{array}
        \right].
    \]
    Eliminating entries below the pivot \(3\), and then eliminating the last
    entry below the third pivot, gives
    \[
        \left[
            \begin{array}{rrrr|r}
                1 & 2 & -1 & 1 & 6\\
                0 & 3 & 1 & 0 & 9\\
                0 & 0 & \frac{17}{3} & 0 & 17\\
                0 & 0 & 0 & 1 & 4
            \end{array}
        \right].
    \]
    Back substitution gives
    \[
        z=4,
        \qquad
        y=3,
        \qquad
        x=2,
        \qquad
        w=1.
    \]
    Therefore the solution is
    \[
        (w,x,y,z)=(1,2,3,4).
    \]
\end{snippetexample}

\end{document}
